With the \OHPC{} repository enabled, we can now begin adding desired components onto the
{\em master} server. This repository provides a number of aliases that group
logical components together in order to help aid in this process. For
reference, a complete list of available group aliases and RPM packages available
via \OHPC{} are provided in Appendix~\ref{appendix:manifest}. To add
support for provisioning services, the following command adds a common base
package followed along with the Warewulf provisioning system. Then the main
Warewulf configuration files are edited to reflect the environment.

%\nottoggle{isCentOS}{\clearpage}

% begin_ohpc_run
% ohpc_comment_header Add baseline OpenHPC\ref{sec:add_provisioning}
\begin{lstlisting}[language=bash,keywords={}]
# Make sure required variables are set
[sms](*\#*) : ${sms_eth_internal:?ERROR: sms_eth_internal not set}

# Install base packages
[sms](*\#*) (*\install*) ohpc-base yq

# Install Warewulf
[sms](*\#*) (*\install*) warewulf-ohpc

# Create the tftpboot directory (not created by dnsmasq)
[sms](*\#*) install -d -m 0755 /var/lib/tftpboot

# Add public_content_t type to tftpboot directory and contents
[sms](*\#*) semanage fcontext -a -t public_content_t "/var/lib/tftpboot(/.*)?"
[sms](*\#*) restorecon -R -v /var/lib/tftpboot

\end{lstlisting}
% end_ohpc_run
